\SourceLeaf{251}%
% A544-P0251-D01
beginnt mit Art.\ 367 und sollte also die Fortsetzung der Disqq.\ Arithm.\ bilden; die wenigen noch vorhandenen Artikel sind aber später ihrem Inhalte nach in die Abhandlung Summatio quarundam serierum singularium übergegangen, und deshalb wird dieses Fragment von der gegenwärtigen Ausgabe ausgeschlossen.

% A544-P0251-D02
In dem vorstehenden Abdruck der beiden Theile der Analysis Residuorum ist der Text des Originals im Wesentlichen treu beibehalten, obgleich dasselbe in formeller Beziehung nicht druckfertig zu nennen ist; in den folgenden Bemerkungen sind die wichtigsten Abänderungen bezeichnet, und zugleich einige Erläuterungen hinzugefügt.

% A544-P0251-D03
\textsection.\ 237.\ Vergl.\ Disqq.\ Arithm.\ artt.\ 61, 62.

% A544-P0251-D04
\textsection.\ 239.\ Vergl.\ Disqq.\ Arithm.\ artt.\ 53, 54, 65.

% A544-P0251-D05
\textsection.\ 241.\ Wenn $n=2^\nu$ und $\nu\geqq3$ ist, so existirt zwar keine Zahl $\rho$ von der angegebenen Art, aber die ganze Untersuchung wird hierdurch nicht wesentlich geändert.

% A544-P0251-D06
\textsection.\ 251.\ Vermuthlich sollte die hier bemerkte Schwierigkeit durch die Einführung höherer Potenzen von $p$ als Moduln beseitigt werden. Vergl.\ \textsection\textsection.\ 363, 372, 373.

% A544-P0251-D07
\textsection.\ 332.\ Die Voraussetzung, dass der Modulus eine Primzahl ist, wird bis \textsection.\ 372 incl.\ beibehalten.

% A544-P0251-D08
\textsection.\ 338.\ Das unvollständige Citat kann auf Disqq.\ Arithm.\ art.\ 44 bezogen werden.

% A544-P0251-D09
\textsection\textsection.\ 344--346.\ Von den beiden im Manuscript vorhandenen Beweisen ist hier der erste, welcher mit den Worten iam demonstrare accingimur eingeleitet wird und sich auf eine nähere Untersuchung der Ausdrücke $(1^\alpha2^\beta3^\gamma\ldots)$ gründet, nach der eigenen Vorschrift des Verfassers ganz unterdrückt (`Tota praecedens demonstratio una cum altera theorematis praec., quam adiicere mens erat, supprimenda erit, quoniam aliam infinities simpliciorem deteximus. Nititur ea huic fundamento'); in dem obigen Abdruck ist ferner der zweite Beweis dadurch abgekürzt, dass die Entwicklung von $\frac{x\,\mathrm{d}P}{P\,\mathrm{d}x}$ statt derjenigen von $\frac{x\,\mathrm{d}P}{\mathrm{d}x}$ betrachtet wird, wodurch zugleich eine im Original enthaltene Beziehung auf den unterdrückten ersten Beweis umgangen wird.

% A544-P0251-D10
\textsection.\ 348.\ Der Ausdruck radix prima ist hier in derselben Bedeutung zu nehmen, wie der Ausdruck radix propria in der Abhandlung Summatio quarundam serierum singularum art.\ 11. --- Bei der Behauptung, dass die Coëfficienten $A'.B'\ldots$ des entwickelten Productes ganze rationale Zahlen sind, wird auf das sechste Capitel verwiesen, in welchem aber die Theorie der Gleichung $x^\tau-1=0$ nur für den Fall behandelt wird, dass $\tau$ eine Primzahl ist; die Form des Beweises in \textsection.\ 349 führt zunächst auf folgende Ergänzung. Wird das entwickelte Product in die (für alle Wurzeln der Gleichung $\theta^\tau=1$ geltende) Form
\[
S=E+F\theta+\ldots+N\theta^{\tau-1}
\]
gebracht, so sind die Coëfficienten $E,F\ldots N$ ganze rationale Functionen von $x$ mit ganzen rationalen Coëfficienten; da ferner das Product ungeändert bleibt, wenn $\theta$ durch $\theta^k$ ersetzt wird, wo $k$ irgend eine relative Primzahl zu $\tau$ bedeutet, so gilt dasselbe von dem Ausdruck $S$, und hieraus ergiebt sich ohne Schwierigkeit, dass alle diejenigen in $S$ enthaltenen Potenzen von $\theta$, deren Exponenten $s$ einen und denselben grössten gemeinschaftlichen Divisor mit $\tau$ haben, auch identische Coëfficienten haben müssen; da endlich eine jede Summe solcher Potenzen $\theta^s$ immer eine ganze Zahl ist, so leuchtet ein, dass der Ausdruck $S$, und folglich auch das in Rede stehende Product eine ganze Function von $x$ mit ganzen Coëfficienten ist, was zu zeigen war. Ebenso geht aus dieser Betrachtung zugleich die Richtigkeit der Bemerkung am Schlusse des Paragraphen hervor. Andere Gründe lassen indessen vermuthen, dass dem Verfasser schon damals das allgemeine Theorem über die Transformation der symmetrischen Functionen (Demonstratio nova altera theorematis omnem functionem etc.\ art.\ 4) bekannt war, aus welchem sich die obigen Sätze als unmittelbare Folgerungen ergeben.

% A544-P0251-D11
\textsection.\ 352.\ Das Zeichen $R\equiv S(\mathrm{mod.}\,P)$ oder auch $R\equiv S(\mathrm{mod.}\,P,p)$ bedeutet hier und im 
\SourceLeaf{252}%
% A544-P0252-D01
Folgenden, dass die Differenz $R-S$ nach dem Modul $p$ den Divisor $P$ hat. --- Das unvollständige Citat kann auf Disqq.\ Arithm.\ art.\ 49 bezogen werden.

% A544-P0252-D02
\textsection.\ 354.\ Durch Multiplication mit $x^\nu-1$ ergiebt sich, dass die Summen gleich hoher Potenzen der Wurzeln der beiden Gleichungen $(P,\rho^{k\nu+t})=0$, $(P,\rho^t)=0$ einander congruent sind $(\mathrm{mod.}\,p)$, und hieraus folgt die Congruenz $(P,\rho^{k\nu+t})\equiv(P,\rho^t)(\mathrm{mod.}\,p)$, sobald $m<p$ ist (vergl.\ \textsection.\ 244); ist aber $m\geqq p$, so lässt sich der Coëfficient der Potenz $x^{m-p}$ in einer Gleichung nicht mehr aus den gegebenen Potenzsummen ihrer Wurzeln nach dem Modul $p$ bestimmen, weil er in den hierzu dienenden \textsc{Newton}'schen Formeln mit dem Factor $p$ behaftet ist. In der That darf man aus der Congruenz je zweier gleich hoher Potenzsummen der Wurzeln der Gleichungen $A=0$, $B=0$ allgemein nur folgern, dass $A\equiv\mathfrak A^p\mathfrak C$, $B=\mathfrak B^p\mathfrak C(\mathrm{mod.}\,p)$ ist, wo $\mathfrak C$ den grössten gemeinschaftlichen Divisor der beiden Functionen $A,B$ nach dem Primzahl-Modulus $p$ bezeichnet, $\mathfrak A$ und $\mathfrak B$ aber ganz unbestimmte Functionen sind. Es ist zu vermuthen, dass der Verfasser die Allgemeingültigkeit des Satzes aus der Theorie der Transformation der symmetrischen Functionen und speciell aus dem folgenden Satze abgeleitet hat: Ist in Bezug auf einen beliebigen Modulus $p$ die Differenz $R(x)-S(x)$ theilbar durch die Function $P(x)$, und sind $a,b,c\ldots$ die Wurzeln der Gleichung $P(x)=0$, so sind die Functionen
\[
\begin{gathered}
(x-R(a))(x-R(b))(x-R(c))\ldots\quad\text{und}\\
(x-S(a))(x-S(b))(x-S(c))\ldots
\end{gathered}
\]
einander nach dem Modul $p$ congruent.

% A544-P0252-D03
\textsection.\ 355.\ Es wird in \textsection.\ 368 gezeigt, dass $P$ und $\frac{\mathrm{d}P}{\mathrm{d}x}$ keinen gemeinschaftlichen Divisor haben, wenn $P$ keinen Factor mehr als einmal enthält.

% A544-P0252-D04
\textsection\textsection.\ 358, 359.\ Die unter den Text gesetzte Note ist einem einzelnen Blatt entnommen, welches wahrscheinlich den schon in der Handschrift gestrichenen \textsection.\ 359 ersetzen sollte.

% A544-P0252-D05
\textsection.\ 360.\ In dem Ausdruck des Theorems ist eine Ungenauigkeit der Handschrift berichtigt.

% A544-P0252-D06
\textsection.\ 361.\ Hier bedeutet der Exponent $\frac1k$ in dem Zeichen $(\xi,\rho^{\frac1k})$ jede positive ganze Zahl $k'$ von der Beschaffenheit, dass $kk'\equiv1(\mathrm{mod.}\,\nu)$ wird, wo $\nu$ die kleinste positive ganze Zahl ist, für welche $x^\nu-1$ durch $\xi$ nach dem Modul $p$ theilbar wird; hierbei ist vorauszusetzen, dass $\xi$ nicht durch $x$ theilbar nach dem Modul $p$, und ausserdem, dass $k$ relative Primzahl zu $\nu$ ist. Die Richtigkeit der Behauptung, dass $\xi'$ durch $(\xi,\rho^{\frac1k})$ theilbar ist $(\mathrm{mod.}\,p)$, ergiebt sich aus \textsection.\ 354.

% A544-P0252-D07
\textsection.\ 363.\ Die Schlussbemerkung bezieht sich vermuthlich auf die Einführung von Moduln, welche Potenzen der Primzahl $p$ sind; vergl.\ \textsection\textsection.\ 251, 372, 373.

% A544-P0252-D08
\textsection.\ 367.\ Die Wurzeln der Gleichung $x^3+xx-2x-1=0$ sind die zweigliedrigen Perioden, in welche die Wurzeln der Gleichung $\frac{x^7-1}{x-1}=0$ zerfallen. Dasselbe Beispiel findet sich auch auf einem einzelnen Blatt, wo das Hauptresultat der \textsection\textsection.\ 362, 363 unter dem Titel `der goldene Lehrsatz' ausgesprochen ist.

% A544-P0252-D09
\textsection.\ 371.\ Dieser Paragraph sollte ein Beispiel enthalten; doch ist dasselbe nicht ausgeführt.
\begin{flushright}R.\ \textsc{Dedekind}.\end{flushright}
\BookRule{.28\linewidth}

\SourceLeaf{253}%
% A544-P0253-D01
\par\bigskip\Needspace{9\baselineskip}
\pdfbookmark[0]{Disquisitionum circa aequationes puras ulterior evolutio}{aequationes-purae}
\begin{center}\large\bfseries DISQUISITIONUM CIRCA AEQUATIONES PURAS\par\medskip ULTERIOR EVOLUTIO.\end{center}
\BookRule{.12\linewidth}
\renewcommand{\BookRunningTitle}{Nachlass. Disquisitionum circa aequationes puras}
\renewcommand{\Article}[1]{\par\medskip\Needspace{4\baselineskip}\pdfbookmark[1]{Art. #1}{purae-art-#1}\centerline{\textbf{#1.}}\nobreak\smallskip}

% A544-P0253-D02
\Article{1}
Quum methodus ea, per quam in \emph{Disquiss.\ Arithm.}\ art.\ 360 aequationem $x^n-1=0$ solvere docuimus, theoriam foecundissimam et gravissimam constituat, cuius prima tantum momenta in opere illo attingere licuit, gratum geometris fore speramus, si hoc argumentum denuo hic resumimus, quae breviter tantum partimque demonstrationibus suppressis adumbrata fuerant, uberius tractamus, et quae ex illo tempore accesserunt incrementa profundius persequimur.

% A544-P0253-D03
Exponens $n$ supponitur esse numerus primus, numerusque $n-1$ in factores $\alpha\times\beta\times\gamma$ resolutus; porro designamus per $g$ aliquam radicem primitivam pro modulo $n$. Exhibeat $r$ indefinite radicem aequationis $x^n-1=0$, atque $R$ indefinite radicem aequationis $x^\beta-1=0$. Designando itaque peripheriam circuli, cuius radius $=1$, per $P$, quantitatemque imaginariam $\mathord{\surd}{-1}$ per $i$, omnes radices aequationis $x^\beta-1=0$, sive omnes valores ipsius $R$ exhibebuntur per formulam
\[
\cos\frac{kP}{\theta}+i\sin\frac{kP}{\theta}
\]
exprimente $k$ indefinite numeros integros $0,1,2,3\ldots\beta-1$. Porro patet, omnes potestates cuiusvis radicis $R$ ipsas quoque esse radices, nec non, si $R$ fuerit radix valori ipsius $k$ ad $\beta$ primo respondens, omnes potestates $R^0,R,R^2,R^3\ldots R^{\beta-1}$ inter se diversas esse, adeoque totum radicum complexum exhaurire; in hoc casu ipsam $R$ radicem \emph{propriam} aequationis $x^\beta-1=0$ dicemus; contra radix $R$ 
\SourceLeaf{254}%
% A544-P0254-D01
valori ipsius $k$ ad $\beta$ non primo respondens \emph{impropria} vocabitur, nulloque negotio perspicitur, si $\delta$ fuerit divisor communis maximus numerorum $k$ et $\beta$, fore $R^{\frac\beta\delta .}=1$, omnes vero potestates $R^0,R,R^2,R^3\ldots R^{\frac\beta\delta-1}$ inter se diversas, adeoque $R$ radicem propriam aequationis $x^{\frac\beta\delta}-1=0$. Eadem de aequatione $x^n-1=0$ valebunt, sed huius radices omnes necessario sunt propriae radice $1$ excepta.

% A544-P0254-D02
\Article{2}
His praemissis disquisitio nostra imprimis versabitur circa functiones huius formae, e $\beta\gamma$ terminis conflatas
\[
r+Rr^{g^\alpha}+R^2r^{g^{2\alpha}}+R^3r^{g^{3\alpha}}\ldots+R^{\beta\gamma-1}r^{g^{\alpha\beta\gamma-\alpha}}
\]
quas compendii caussa per hunc characterem $[r,R]$ designabimus. Singuli termini talis expressionis sunt producta e potestatibus ipsius $r$ in potestates ipsius $R$; illarum exponentes progressionem geometricam constituunt, exponentes harum arithmeticam. Exponentes
\[
1,\quad g^\alpha,\quad g^{2\alpha},\quad g^{3\alpha}\ldots\ldots g^{\alpha\beta\gamma-\alpha}
\]
omnes inter se incongrui sunt secundum modulum $n$, adeoque illae potestates ipsius $r$ inter se diversae; ulterius vero continuatae eandem seriem denuo inciperent, quum sit $g^{\alpha\beta\gamma}\equiv1(\mathrm{mod.}\,n)$, adeoque $r^{g^{\alpha\beta\gamma}}=r$. Factores alteri autem
\[
1,\quad R,\quad R^2,\quad R^3\dotsfour R^{\beta\gamma-1}
\]
constituunt $\gamma$ periodos aequales, quum sit $R^\beta=1$, $R^{\beta-1}=R$ etc. Hinc patet, functionem $[r,R]$ ita quoque exhiberi posse
\[
\begin{aligned}
&r+r^{g^{\alpha\beta}}+r^{g^{2\alpha\beta}}\ldots\ldots+r^{g^{\alpha\beta\gamma-\alpha\beta}}\\
&{}+R\bigl(r^{g^\alpha}+r^{g^{\alpha\beta+\alpha}}+r^{g^{2\alpha\beta+\alpha}}\ldots\ldots+r^{g^{\alpha\beta\gamma-\alpha\beta+\alpha}}\bigr)\\
&{}+R^2\bigl(r^{g^{2\alpha}}+r^{g^{\alpha\beta+2\alpha}}+r^{g^{2\alpha\beta+2\alpha}}\ldots\ldots+r^{g^{\alpha\beta\gamma-\alpha\beta+2\alpha}}\bigr)\\
&{}+\text{etc.}\\
&{}+R^{\beta-1}\bigl(r^{g^{\alpha\beta-\alpha}}+r^{g^{2\alpha\beta-\alpha}}+r^{g^{3\alpha\beta-\alpha}}\ldots\ldots+r^{g^{\alpha\beta\gamma-\alpha}}\bigr).
\end{aligned}
\]
sive introducendo signum art.\ 343 Disq.\ Ar.
\[
[r,R]=(\gamma,1)+R(\gamma,g^\alpha)+R^2(\gamma,g^{2\alpha})\ldots+R^{\beta-1}(\gamma,g^{\alpha\beta-\alpha})
\]

\SourceLeaf{255}%
% A544-P0255-D01
\Article{3}
Si pro radice $r$ unitatem accipimus, habemus
\[
[1,R]=1+R+R^2+R^3\ldots+R^{\beta\gamma-1}=\gamma(1+R+R^2+R^3\ldots+R^{\beta-1})
\]
huius valor erit $=\beta\gamma$, si etiam pro $R$ accipitur radix $1$, sed $=0$ pro quovis alio valore ipsius $R$. Contra manente $r$ indeterminata, positaque $R=1$, erit $[r,1]=r+r^{g^\alpha}+r^{g^{2\alpha}}+r^{g^{3\alpha}}\ldots+r^{g^{\alpha\beta\gamma-\alpha}}$, sive adhibito signo in Disq.\ Ar.\ introducto, $[r,1]=(\beta\gamma,1)$, i.\ e.\ constabit e periodo $\beta\gamma$ radicum, e quibus una est ipsa $r$. Quoties est $\alpha=1$, haec periodus omnes radices $r,r^2,r^3\ldots r^{n-1}$ complectetur ordine tantum mutato.

% A544-P0255-D02
Notentur adhuc relationes sequentes, quarum ratio sponte elucet:
\[
[r,R]=R[r^{g^\alpha},R]=R^2[r^{g^{2\alpha}},R]\quad\text{sive generaliter}\quad=R^k[r^{g^{\alpha k}},R]
\]
denotante $k$ integrum positivum quemcunque. Hinc patet, functionem $[r^m,R]$ vel esse $=[1,R]$, scilicet si fuerit $m$ divisibilis per $n$, vel reduci posse ad formam $R^\mu[r^{g^\nu},R]$ in casibus reliquis et quidem ita, ut sit $\nu<\alpha$. Si enim $m$ non est divisibilis per $n$, congruus erit secundum modulum $n$ alicui potestati ipsius $g$, cuius exponens ad instar Disq.\ Ar.\ per $\mathrm{ind.}\,m$ commode exprimitur; statuendo itaque $\mathrm{ind.}\,m=\lambda\alpha+\nu$, quod manifesto fieri potest, ita ut sit $\nu<\alpha$, erit $[r^m,R]=[r^{g^{\lambda\alpha+\nu}},R]=R^{-\lambda}[r^{g^\nu},R]$: faciendus est itaque $\mu=-\lambda$ aut si exponentem positivum desideras, $\mu\equiv-\lambda(\mathrm{mod.}\,\beta)$.

% A544-P0255-D03
\Article{4}
\textsc{Theorema.} \emph{Designante $r'$ perinde ut $r$ indefinite radicem aequationis $x^n-1=0$, nec non $R'$ perinde ut $R$ indefinite radicem aequationis $x^\beta-1=0$, erit productum}
\[
\begin{aligned}
&[r,R]\times[r',R']=\\
&\quad[rr',RR']+R[r^{g^\alpha}r',RR']+R^2[r^{g^{2\alpha}}r',RR']\\
&\quad{}+R^3[r^{g^{3\alpha}}r',RR']\ldots+R^{\beta\gamma-1}[r^{g^{\alpha\beta\gamma-\alpha}}r',RR']
\end{aligned}
\]

% A544-P0255-D04
\emph{Demonstr.} Absolvendo multiplicationem ipsius $[r,R]$ per singulas partes ipsius $[r',R']$, productum in hac forma exhiberi potest
\[
\begin{aligned}
&[r,R]r'+RR'[r^{g^\alpha},R]r'^{g^\alpha}+R^2R'^2[r^{g^{2\alpha}},R]r'^{g^{2\alpha}}\\
&{}+R^3R'^3[r^{g^{3\alpha}},R]r'^{g^{3\alpha}}\ldots\\
&{}+R^{\beta\gamma-1}R'^{\beta\gamma-1}[r^{g^{\alpha\beta\gamma-\alpha}},R]r'^{g^{\alpha\beta\gamma-\alpha}}
\end{aligned}
\]

